% Regression: t2_eq54_ginj — renders (1/|G|) sum_{g in G} L_g (x) L_g (x) L_g^{-1} (x) L_g^{-1} [one action per…
% T2 — RMP arXiv:2011.12127, fig. fig3-pepe_Eq54Ginjred (eq. 54,
% G-injectivity): the group average of the on-bond symmetry action equals
% the red symmetry ring threaded through the plaquette's virtual indices.
% Formula:
%   (1/|G|) sum_{g in G}  L_g (x) L_g (x) L_g^{-1} (x) L_g^{-1}
%     [one action per virtual leg of the plaquette, arrows = V vs V*]
%   =  P_ring : the projector drawn as the red virtual ring passing
%      through the four surrounding PEPS tensors, with the orientation
%      dot on its NW corner.
% Free tier (genuinely non-grid: a cross of directed legs, then a closed
% ring with tensors on all four EDGES -- no grid row can hold N/E/S/W
% placements).  L_g = the on-wire-matrix skin (`ring`); the ring = four
% hv/vh joins with the orientation dot AT the NW corner.
%
% Missing primitive (checked in tenkz-free.code.tex, \tnjoin style
% resolution): joins stroke `bond`/`fused bond` only -- there is no
% operator-hued join (no /tenkz/join/operator key), so the symmetry ring
% that the source draws RED renders black here.  The hue is load-bearing
% (it separates the symmetry string from the black lattice legs).
\documentclass[varwidth=19cm, border=6pt]{standalone}
\usepackage{amsmath, amssymb}
\usepackage{tenkz}
\begin{document}
\[
  \frac{1}{|G|}\sum_{g\in G}\;
  % LHS: one group action per virtual leg; barbs give V (out) vs V* (in)
  \begin{tenkzfree}
    \tnput[ring]{lw}{(-0.8,0)}{L_g}
    \tnput[ring]{le}{(0.8,0)}{L_g}
    \tnput[ring]{lt}{(0,0.8)}{L_g}
    \tnput[ring]{lb}{(0,-0.8)}{L_g}
    % west leg flows east, east leg flows west (the inverse action),
    % top flows down, bottom flows up: all toward the plaquette centre
    \tnjoin[dir]{-1.45,0}{lw.west}   \tnjoin[dir]{lw.east}{-0.2,0}
    \tnjoin[dir]{1.45,0}{le.east}    \tnjoin[dir]{le.west}{0.2,0}
    \tnjoin[dir]{0,1.45}{lt.north}   \tnjoin[dir]{lt.south}{0,0.2}
    \tnjoin[dir]{0,-1.45}{lb.south}  \tnjoin[dir]{lb.north}{0,-0.2}
  \end{tenkzfree}
  \;=\;
  % RHS: the red ring through the four PEPS tensors, dot on the NW corner
  \begin{tenkzfree}[tensor style=box]
    \tnput[box]{n}{(0,0.8)}{}
    \tnput[box]{e}{(0.8,0)}{}
    \tnput[box]{s}{(0,-0.8)}{}
    \tnput[box]{w}{(-0.8,0)}{}
    \tnput[dot]{d}{(-0.8,0.8)}{}
    % black lattice legs through each tensor (physical ink of the source)
    \tnjoin{-1.45,0}{w.west}  \tnjoin{w.east}{-0.35,0}
    \tnjoin{1.45,0}{e.east}   \tnjoin{e.west}{0.35,0}
    \tnjoin{0,1.45}{n.north}  \tnjoin{n.south}{0,0.35}
    \tnjoin{0,-1.45}{s.south} \tnjoin{s.north}{0,-0.35}
    % the symmetry ring (source: RED) with the orientation dot at NW
    \tnjoin{d.east}{n.west}
    \tnjoin[route=hv]{n.east}{e.north}
    \tnjoin[route=vh]{e.south}{s.east}
    \tnjoin[route=hv]{s.west}{w.south}
    \tnjoin{w.north}{d.south}
  \end{tenkzfree}
\]
\end{document}
